\documentclass{article}
\usepackage[utf8]{inputenc}
\usepackage[ngerman]{babel}
\usepackage{booktabs}
\usepackage{amsmath}
\usepackage{tabularx}

\begin{document}
	
	\section*{Systematik des Periodensystems der Elemente ($Z = 1 \dots 93$)}
	
	\subsection*{1. Periode 1 bis 3 (Leichte Elemente)}
	\begin{table}[htbp]
		\centering
		\small
		\begin{tabular}{ccll}
			\toprule
			$\mathbf{Z}$ & \textbf{Symbol} & \textbf{Element} & \textbf{Kategorie / Block} \\
			\midrule
			1  & H  & Wasserstoff & Nichtmetall ($s$-Block) \\
			2  & He & Helium      & Edelgas ($s$-Block) \\
			3  & Li & Lithium     & Alkalimetall \\
			4  & Be & Beryllium   & Erdalkalimetall \\
			5  & B  & Bor         & Halbmetall \\
			6  & C  & Kohlenstoff & Nichtmetall \\
			7  & N  & Stickstoff  & Nichtmetall \\
			8  & O  & Sauerstoff  & Nichtmetall \\
			9  & F  & Fluor       & Halogen \\
			10 & Ne & Neon        & Edelgas \\
			11 & Na & Natrium     & Alkalimetall \\
			12 & Mg & Magnesium   & Erdalkalimetall \\
			13 & Al & Aluminium   & Metall \\
			14 & Si & Silicium    & Halbmetall \\
			15 & P  & Phosphor    & Nichtmetall \\
			16 & S  & Schwefel    & Nichtmetall \\
			17 & Cl & Chlor       & Halogen \\
			18 & Ar & Argon       & Edelgas \\
			\bottomrule
		\end{tabular}
	\end{table}
	
	\subsection*{2. Periode 4 (Erste Nebengruppe / $d$-Block)}
	\begin{table}[htbp]
		\centering
		\small
		\begin{tabular}{cclccl}
			\toprule
			$\mathbf{Z}$ & \textbf{Symbol} & \textbf{Element} & $\mathbf{Z}$ & \textbf{Symbol} & \textbf{Element} \\
			\midrule
			19 & K  & Kalium    & 28 & Ni & Nickel \\
			20 & Ca & Calcium   & 29 & Cu & Kupfer \\
			21 & Sc & Scandium  & 30 & Zn & Zink \\
			22 & Ti & Titan     & 31 & Ga & Gallium \\
			23 & V  & Vanadium  & 32 & Ge & Germanium \\
			24 & Cr & Chrom     & 33 & As & Arsen \\
			25 & Mn & Mangan    & 34 & Se & Selen \\
			26 & Fe & Eisen     & 35 & Br & Brom \\
			27 & Co & Cobalt    & 36 & Kr & Krypton \\
			\bottomrule
		\end{tabular}
	\end{table}
	
	\subsection*{3. Periode 5}
	\begin{table}[htbp]
		\centering
		\small
		\begin{tabular}{cclccl}
			\toprule
			$\mathbf{Z}$ & \textbf{Symbol} & \textbf{Element} & $\mathbf{Z}$ & \textbf{Symbol} & \textbf{Element} \\
			\midrule
			37 & Rb & Rubidium   & 46 & Pd & Palladium \\
			38 & Sr & Strontium  & 47 & Ag & Silber \\
			39 & Y  & Yttrium    & 48 & Cd & Cadmium \\
			40 & Zr & Zirkonium  & 49 & In & Indium \\
			41 & Nb & Niob       & 50 & Sn & Zinn \\
			42 & Mo & Molybdän   & 51 & Sb & Antimon \\
			43 & Tc & Technetium & 52 & Te & Tellur \\
			44 & Ru & Ruthenium  & 53 & I  & Iod \\
			45 & Rh & Rhodium    & 54 & Xe & Xenon \\
			\bottomrule
		\end{tabular}
	\end{table}
	
	\subsection*{4. Periode 6 (Inklusive Lanthanoide)}
	\begin{table}[htbp]
		\centering
		\small
		\caption{Lanthanoide ($f$-Block, $Z = 57 \dots 71$)}
		\begin{tabular}{cclccl}
			\toprule
			$\mathbf{Z}$ & \textbf{Symbol} & \textbf{Element} & $\mathbf{Z}$ & \textbf{Symbol} & \textbf{Element} \\
			\midrule
			57 & La & Lanthan     & 65 & Tb & Terbium \\
			58 & Ce & Cer         & 66 & Dy & Dysprosium \\
			59 & Pr & Praseodym   & 67 & Ho & Holmium \\
			60 & Nd & Neodym      & 68 & Er & Erbium \\
			61 & Pm & Promethium  & 69 & Tm & Thulium \\
			62 & Sm & Samarium    & 70 & Yb & Ytterbium \\
			63 & Eu & Europium    & 71 & Lu & Lutetium \\
			64 & Gd & Gadolinium  &    &    & \\
			\bottomrule
		\end{tabular}
	\end{table}
	
	\subsection*{5. Periode 7 (Bis zum ersten Transuran)}
	\begin{table}[htbp]
		\centering
		\small
		\begin{tabular}{ccll}
			\toprule
			$\mathbf{Z}$ & \textbf{Symbol} & \textbf{Element} & \textbf{Status / Anmerkung} \\
			\midrule
			87 & Fr & Francium    & Alkalimetall \\
			88 & Ra & Radium      & Erdalkalimetall \\
			89 & Ac & Actinium    & Beginn der Actinoide \\
			90 & Th & Thorium     & Actinoid (natürlich vorkommend) \\
			91 & Pa & Protactinium & Actinoid (natürlich vorkommend) \\
			92 & U  & Uran        & Schwerstes natürlich vorkommendes Element \\
			\midrule
			93 & Np & Neptunium   & \textbf{Erstes Transuran} ($Z > 92$, künstlich erzeugt) \\
			\bottomrule
		\end{tabular}
	\end{table}
	
\end{document}